%-------------------------
% Based off of: https://github.com/sb2nov/resume
% License : MIT
%------------------------

\documentclass[letterpaper,10pt]{article}

\usepackage{latexsym}
\usepackage[empty]{fullpage}
\usepackage{titlesec}
\usepackage{marvosym}
\usepackage[usenames,dvipsnames]{color}
\usepackage{verbatim}
\usepackage{enumitem}
\usepackage[hidelinks]{hyperref}
\usepackage{fancyhdr}
\usepackage[russian, english]{babel}
\usepackage{tabularx}
\input{glyphtounicode}


%----------FONT OPTIONS----------
% sans-serif
% \usepackage[sfdefault]{FiraSans}
% \usepackage[sfdefault]{roboto}
% \usepackage[sfdefault]{noto-sans}
% \usepackage[default]{sourcesanspro}

% serif
% \usepackage{CormorantGaramond}
% \usepackage{charter}


\pagestyle{fancy}
\fancyhf{} % clear all header and footer fields
\fancyfoot{}
\renewcommand{\headrulewidth}{0pt}
\renewcommand{\footrulewidth}{0pt}

% Adjust margins
\addtolength{\oddsidemargin}{-0.5in}
\addtolength{\evensidemargin}{-0.5in}
\addtolength{\textwidth}{1in}
\addtolength{\topmargin}{-.5in}
\addtolength{\textheight}{1.0in}

\urlstyle{same}

\raggedbottom
\raggedright
\setlength{\tabcolsep}{0in}

% Sections formatting
\titleformat{\section}{
	\vspace{-4pt}\scshape\raggedright\large
}{}{0em}{}[\color{black}\titlerule \vspace{-5pt}]

% Ensure that generate pdf is machine readable/ATS parsable
\pdfgentounicode=1

%-------------------------
% Custom commands
\newcommand{\resumeItem}[2]{
	\item\small{
		{\ifx&#1&\else\textbf{#1}: \fi#2 \vspace{-2pt}}
	}
}

\usepackage[T2A]{fontenc}
\usepackage[utf8]{inputenc}
\usepackage{cmap}

\newcommand{\resumeSubheading}[4]{
	\vspace{-1pt}\item
	\begin{tabular*}{0.97\textwidth}{l@{\extracolsep{\fill}}r}
		\textbf{#1} & #2 \\
		\ifx&#3&\else\textit{\small#3} & \textit{\small #4} \\ \fi
	\end{tabular*}\vspace{-5pt}
}


\newcommand{\resumeSubSubheading}[2]{
	\item
	\begin{tabular*}{0.97\textwidth}{l@{\extracolsep{\fill}}r}
		\textit{\small#1} & \textit{\small #2} \\
	\end{tabular*}\vspace{-7pt}
}

\newcommand{\resumeProjectHeading}[2]{
	\item
	\begin{tabular*}{0.97\textwidth}{l@{\extracolsep{\fill}}r}
		\small#1 & #2 \\
	\end{tabular*}\vspace{-7pt}
}

\newcommand{\resumeSubItem}[2]{\resumeItem{#1}{#2}\vspace{-4pt}}

\renewcommand\labelitemii{$\vcenter{\hbox{\tiny$\bullet$}}$}

\newcommand{\resumeSubHeadingListStart}{\begin{itemize}[leftmargin=0.15in, label={}]}
	\newcommand{\resumeSubHeadingListEnd}{\end{itemize}}
\newcommand{\resumeItemListStart}{\begin{itemize}}
	\newcommand{\resumeItemListEnd}{\end{itemize}\vspace{-5pt}}
\newcommand{\ulhref}[2]{\href{#1}{\underline{#2}}}
%-------------------------------------------
%%%%%%  RESUME STARTS HERE  %%%%%%%%%%%%%%%%%%%%%%%%%%%%


\begin{document}

	\begin{center}
		\textbf{\Large \scshape Усатов Павел} \textbf{\Large $|$ Python Developer} \\ \vspace{1pt}
		\small +7 (912) 712-83-52 $|$
		\href{https://t.me/Usatov_Pavel}{\underline{Telegram}} $|$
		\href{mailto:pvusatov@edu.hse.ru}{\underline{pvusatov@edu.hse.ru}} $|$
		\href{https://github.com/UsatovPavel}{\underline{Github}} $|$
		\href{https://usatovpavel.ru}{\underline{\textbf{Сайт}}}
	\end{center}

	%-----------EDUCATION-----------
	\section{Образование}
	\resumeSubHeadingListStart
	\resumeSubheading
	{ВШЭ СПБ}{Санкт-Петербург}
	{Прикладная математика и информатика}{2023-2027}
	\item \textbf{Релевантные курсы:}
	Алгоритмы и структуры данных, базы данных, компьютерные сети,
	архитектура компьютера и операционные системы, архитектура ПО
	\resumeSubHeadingListEnd

	%-----------EXPERIENCE-----------
	\section{Опыт}
	\resumeSubHeadingListStart
	\resumeSubheading
	{Стажировка Backend Т-Банк}{апрель 2026 - июль 2026}
	{Участвовал в миграции документов с long на uuid}{Продуктовая команда документооборота}
	\resumeItemListStart
	\resumeItem{}{Версионировал 10 эндпоинтов под uuid, чтобы не сломать мобильные приложения, установленные у пользователей}
	\resumeItem{}{Написал модуль команды в репозитории интеграции: 7 клиентов взамен legacy, смежные команды мигрируют на API}
	\resumeItem{}{Чинил баги генерации документов: время у нас, восстановил столбец статистики во внешнем сервисе}
	\resumeItem{}{Работа с БД: создание алертов, оптимизация запроса, добавление расширений}
	\resumeItemListEnd
	\textbf{Технологии}: Scala, Cats Effect, PostgreSQL, Kafka, Docker, Typescript, Kanban
	\resumeSubheading
	{\ulhref{https://github.com/UsatovPavel/riid}{Утилита RIID для внутреннего облака VK}}{январь 2026 - июнь 2026}
	{Daemon для p2p-загрузки OCI/Docker образов \textit{независимый} от container engine}{Курсовая, научрук - SRE из VK}
	\resumeItemListStart
	\resumeItem{}{Написал \ulhref{https://github.com/UsatovPavel/java-dragonfly-image-puller}{библиотеку}: gRPC-клиент к Dragonfly для загрузки слоёв с соседних нод. \textbf{Взят в использование в VK}}
	\resumeItem{}{На 20\% быстрее Podman на кластере k8s из 12 нод: 100 образов от 1 МБ до 5 ГБ}
	\resumeItemListEnd
	\textbf{Технологии}: Java, Kubernetes, gRPC, Docker, Podman, OCI/Docker Registry API, p2p, Grafana

	\resumeSubHeadingListEnd

	%-----------PROJECTS-----------
	\section{Проекты}
	\resumeSubHeadingListStart
	\resumeSubheading
	{\ulhref{https://github.com/UsatovPavel/Pop_ML_Forecast10}{\textbf{Прогноз населения стран}} $|$ \normalfont{\emph{ML на 10-летнем горизонте}}}{март 2026}
	{Воспроизведение и расширение статьи по ML-прогнозированию демографии}{Индивидуальный проект}
	\resumeItemListStart
	\resumeItem{}{Собрал датасет из UN WPP, US Census IDB и World Bank API: когорты, смертность, фертильность}
	\resumeItem{}{Сравнил 10 конфигураций: ARIMA с grid search, Prophet, XGBoost, Random Forest, когортная модель}
	\resumeItem{}{Лучший MAPE 0.020 против 0.196 у наивного baseline; прогноз по России сверен с росстатовским}
	\resumeItemListEnd
	\textbf{Технологии}: Python, pandas, NumPy, scikit-learn, statsmodels, XGBoost, Prophet, Jupyter

	\resumeSubheading
	{\ulhref{https://github.com/FUYOH666/scanovich-webUI/}{\textbf{Хакатон MTS True Tech}} $|$ \normalfont{\emph{оркестратор нейросетей}}}{апрель 2026}
	{Приватный AI-контур заказчика, интеграция с Open WebUI}{Командный проект}
	\resumeItemListStart
	\resumeItem{}{Генерация PPTX по плану от LLM: шаблоны и адаптация содержания под аудиторию}
	\resumeItem{}{Семантический классификатор запросов, политика ingest, мост статусов Open WebUI, STT}
	\resumeItemListEnd
	\textbf{Технологии}: Python, LLM API, Open WebUI, Docker Compose, Make

	\resumeSubheading
	{\ulhref{https://github.com/hse-project-Java-2025}{TimeTamer}}{2025}
	{Java server$+$Kotlin Android приложение-календарь с элементами AI и геймификации}{Командный проект}
	\resumeItemListStart
	\resumeItem{}{Интеграция AI-ассистента (ChatGPT + Whisper) для предложений по задачам и голосового ввода}
	\resumeItemListEnd
	\textbf{Технологии}: Java, Spring Boot, Kotlin, PostgreSQL, Docker, JUnit 4/5, Mockito

	\resumeSubheading
	{\textbf{pr-review-rag} $|$ \normalfont{\emph{прототип приватного AI-ревью пул-реквестов}}}{июль 2026}
	{RAG по чанкам diff в SQLite, GigaChat и Yandex AI Studio за одним интерфейсом}{Индивидуальный проект}
	\resumeItemListStart
	\resumeItem{}{За 7 часов довёл до рабочего прототипа и закрыл направление: качество ревью неудовлетворительное}
	\resumeItemListEnd
	\textbf{Технологии}: Python, sentence-transformers, SQLite, httpx, GigaChat API, Yandex AI Studio

	\resumeSubheading
	{Учебные задания}{}{}{}
	\resumeItemListStart
	\resumeItem{Python}{\ulhref{https://github.com/UsatovPavel/Usatov-FL-HSE/tree/task4-dev}{ANTLR-парсер} регулярных выражений: грамматика, генерация инструкций, pytest}
	\resumeItem{Go}{\ulhref{https://github.com/UsatovPavel/PRAssign}{REST API на Gin} с JWT, пайплайн Kafka $\rightarrow$ PostgreSQL, нагрузочные тесты k6}
	\resumeItemListEnd
	\resumeSubHeadingListEnd

	%-----------SKILLS-----------
	\section{Навыки}
	\textbf{Языки}{: Python, Java, Scala} \\
	\textbf{Python}{: pandas, NumPy, scikit-learn, statsmodels, XGBoost, Prophet, sentence-transformers, pytest, httpx, ANTLR} \\
	\textbf{Технологии и инструменты}{: Docker, PostgreSQL, SQLite, Kafka, Kubernetes, gRPC, Git, Make, Jupyter} \\
	\textbf{Навыки}{: интеграция LLM и RAG, обработка данных и временные ряды, тестирование (интеграционное/e2e/нагрузочное), алгоритмы и структуры данных}

%-------------------------------------------
\end{document}
